\documentclass[11pt,letterpaper]{article}

\usepackage[margin=1in]{geometry}
\usepackage{amsmath,amssymb}
\usepackage{booktabs}
\usepackage{xcolor}
\usepackage[hidelinks]{hyperref}
\usepackage{microtype}
\usepackage{enumitem}
\usepackage{titlesec}
\usepackage{caption}

\definecolor{mDarkTeal}{HTML}{1c3d6b}
\definecolor{mLightBrown}{HTML}{2f6fb3}

\titleformat{\section}{\large\bfseries\color{mDarkTeal}}{\thesection}{0.6em}{}
\titleformat{\subsection}{\normalsize\bfseries}{\thesubsection}{0.5em}{}
\titlespacing*{\section}{0pt}{1.4em}{0.6em}
\titlespacing*{\subsection}{0pt}{1.0em}{0.4em}

\hypersetup{colorlinks=true, linkcolor=mDarkTeal, citecolor=mDarkTeal, urlcolor=mLightBrown}

\newcommand{\Gop}{\mathcal{G}}
\newcommand{\Hil}{\mathcal{H}}
\newcommand{\D}{\mathrm{D}}

\title{\textbf{CS-DNO: a Craig--Sulem-Structured Neural Dirichlet--Neumann Operator}\\[2pt]
       \large Architecture reference for \texttt{models/dno-net/dno\_net\_v2.py}
       (baseline run \texttt{cs\_dno\_w512b8\_l256\_v3\_20260605\_032608})}
\author{}
\date{June 11, 2026}

\begin{document}
\maketitle

\section{What the model replaces, and where it sits}

The water-wave system is evolved in Zakharov form: with $\eta(x,t)$ the free
surface and $\xi(x,t)$ the surface trace of the velocity potential,
\begin{align}
\eta_t &= \Gop(\eta;h)\,\xi, \\
\xi_t  &= -g\eta \;-\; \tfrac12 \xi_x^2
          \;+\; \frac{\bigl(\Gop(\eta;h)\xi + \eta_x \xi_x\bigr)^2}{2\,(1+\eta_x^2)},
\end{align}
where $\Gop(\eta;h)$ is the \emph{Dirichlet--Neumann operator} (DNO): solve
Laplace's equation in the fluid domain $\{-h < y < \eta(x)\}$ with Dirichlet
data $\xi$ on the surface and no-flux at the flat bottom, and return the
(scaled) normal derivative of the potential at the surface. Every right-hand-side
evaluation requires one application of $\Gop$; the reference implementation is
the order-6 Craig--Sulem series with $8\times$ zero-padding in f64
(\texttt{dno\_series\_eval}), which dominates the cost of a rollout. CS-DNO is a
neural surrogate for the map
\[
(\eta,\ \xi,\ h)\ \longmapsto\ \Gop(\eta;h)\,\xi
\]
on the $n_x=1024$ periodic grid ($L=2\pi$), trained on snapshots harvested from
f64 solver trajectories. At rollout time it is dropped into the same time
integrator in place of the series evaluation (\texttt{solver/evals/model\_rollout.py}).
One convention matters everywhere: the entire data pipeline --- solver RHS,
post-step states, and stored labels --- lives in the dealiased subspace
$k \le 0.25\,k_{\max}$ (cutoff $k=128$), because the order-6 series amplifies
high-$k$ roundoff like $k^6$. The surrogate therefore only ever sees and is only
ever scored on that subspace.

\section{Mathematical structure the architecture encodes}

The DNO admits the Craig--Sulem expansion
\[
\Gop(\eta;h) = \sum_{j\ge 0} \Gop_j(\eta;h),
\qquad \Gop_j \ \text{homogeneous of degree } j \text{ in } \eta,
\]
whose first terms (finite depth, $\D = -i\partial_x$, $|\D| = $ multiplier $|k|$) are
\begin{align}
\Gop_0(h) &= |\D|\tanh(h|\D|), \\
\Gop_1(\eta;h)\,\xi &= -\,\Gop_0\bigl(\eta\, \Gop_0 \xi\bigr)\;-\;\partial_x\bigl(\eta\,\partial_x \xi\bigr).
\end{align}
Each $\Gop_j$ is built by \emph{alternating Fourier multipliers
$m(\D,h)$ with pointwise multiplications by powers of $\eta$}. Four structural
properties of the true operator drive the design:
\begin{enumerate}[leftmargin=1.6em, itemsep=2pt]
\item[\textbf{P1}] \textbf{Linearity in $\xi$.} $\Gop(\eta;h)$ is a linear
      operator applied to $\xi$; all nonlinearity enters through $\eta$.
\item[\textbf{P2}] \textbf{Analyticity in $\eta$.} The expansion converges for
      moderate slopes; the $\eta$-dependence is polynomial order by order.
\item[\textbf{P3}] \textbf{Known dispersion.} The linear part
      $\Gop_0 = |\D|\tanh(h|\D|)$ is exact and closed-form, with all depth
      dependence through $\tanh(hk)$.
\item[\textbf{P4}] \textbf{Multiplier--multiplication factorization.} Higher
      corrections look like $m_{\mathrm{out}}(\D,h)\,[\,\phi(\eta)\cdot
      m_{\mathrm{in}}(\D,h)\,\xi\,]$ composed and summed.
\end{enumerate}
CS-DNO hard-codes P1, P3 and P4 and soft-codes P2. By contrast, a generic FNO
fed $(\eta,\xi)$ as channels respects none of them: it can (and measurably does)
produce output nonlinear in $\xi$, which on near-linear states materializes as
spurious sideband energy; and it must spend capacity re-learning
$\tanh(hk)$ from data. The predecessor \texttt{SpectralDNO} already enforced P1
via a single rank-$r$ ``spectral sandwich'' $C(\eta)^{\!\top} F(\eta)\, C(\eta)$
added to $\Gop_0$, but its $\eta$-trunk was an FNO-style stack of spectral
convolutions and its per-state filter $F$ was produced by a global pooled MLP
--- depth entered only as an input channel. CS-DNO replaces that with a sum of
explicitly Craig--Sulem-shaped blocks and analytic depth conditioning.

\section{Architecture}

\subsection{Interface and normalization}

The module signature is \texttt{(inputs, depth) $\to$ gxi} with
\texttt{inputs} of shape $(B, 1024, 2)$ holding normalized $(\hat\eta, \hat\xi)$
and \texttt{depth} of shape $(B,1)$ holding $\log h$. Under \texttt{norm=scale}
each quantity is divided by its training-dataset absolute maximum:
\[
\hat\eta = \eta / s_\eta, \qquad
\hat\xi  = \xi  / s_\xi, \qquad
\widehat{\Gop\xi} = \Gop\xi / s_g,
\]
with $s_\xi = 0.14096$ and $s_g = 0.76948$ for the v3 dataset (recorded in
\texttt{config.json}; recomputed per dataset). Internally the model un-scales
$\xi$ to physical units before every operator application and re-scales the
output, so the learned parts operate on physically meaningful fields. Depth is
clipped as $\log h \mapsto \min(\log h, \log 5)$: for $k \ge 1$,
$\tanh(5k)$ is within $10^{-4}$ of 1, so all depths beyond $h=5$ are physically
indistinguishable from infinite depth and the clip simply saturates that
direction of input space.

\subsection{Exact linear baseline}

The flat-surface DNO is applied analytically in Fourier space,
\[
\Gop_0(h)\,\xi = \mathrm{irfft}\bigl[\,k\tanh(hk)\cdot \mathrm{rfft}(\xi)\,\bigr],
\]
and the network only ever learns the \emph{residual} $\Gop - \Gop_0$. Because
every learned contribution is zero-initialized (below), the model output at
step~0 of training is exactly $\Gop_0\xi$ --- training starts from the linear
water-wave theory rather than from noise.

\subsection{The $\eta$ feature stack}

The surface shape is expanded into seven fixed analytic channels --- pointwise
powers plus spectral derivatives:
\[
\bigl[\ \eta,\ \eta^2,\ \eta^3,\ \partial_x\eta,\ \partial_x^2\eta,\ |\D|^{1/2}\eta,\ \Hil\eta\ \bigr],
\]
where $\Hil$ is the Hilbert transform ($-i\,\mathrm{sgn}(k)$) and the
half-derivative $|\D|^{1/2}$ is natural in deep-water $\Gop_1$ scalings. The
powers give the polynomial orders of P2 directly; the derivatives expose slope
and curvature (the physical small parameters of the expansion); $\Hil$ and
$|\D|^{1/2}$ give the trunk access to odd/phase-shifted and fractional
information about $\eta$ that pointwise powers cannot produce. These seven
channels pass through a small pointwise trunk shared by all blocks,
\[
\texttt{Dense}(7 \to 256) \to \mathrm{gelu} \to \texttt{Dense}(256 \to 256) \to \mathrm{gelu},
\]
producing $\eta$-features $E(\eta) \in \mathbb{R}^{B \times 1024 \times 256}$
(width$/2$ channels). Note the trunk is \emph{pointwise in $x$} apart from the
fixed spectral derivatives --- there are no learned convolutions over $x$
anywhere in the model; all learned spatial coupling happens through Fourier
multipliers.

\subsection{Depth-aware Fourier multipliers}

Every learned multiplier is produced by \texttt{DepthAwareMultiplier}: a 2-layer
MLP evaluated \emph{per mode} on physical dispersion features,
\[
m(k,h) \;=\; W_2\,\tanh\!\bigl(W_1\,[\,k,\ h,\ \tanh(hk),\ k\tanh(hk)\,] + b_1\bigr) + b_2
\;\in\; \mathbb{R}^{B \times 513 \times r},
\]
with hidden width 128 and $r$ output channels (one per branch). Because the MLP
sees $\tanh(hk)$ and $k\tanh(hk)$ as inputs, the exact dispersion scaling --- and
in particular the correct shallow/deep crossover in $kh$ --- is available as a
\emph{linear readout} rather than something to be discovered; the network learns
perturbations on physics, not the physics itself. This is the entire mechanism
of depth conditioning: $h$ never enters as a spatial channel, only through these
multipliers (and the $\Gop_0$ baseline). The multipliers are real-valued, i.e.
zero-phase filters (a representational choice discussed in \S\ref{sec:expressivity}).

\subsection{The Craig--Sulem block}

Each of the $n_{\mathrm{blocks}}=8$ blocks computes a sum of $r=256$ branches
shaped exactly like a Craig--Sulem term (P4):
\[
\mathrm{Block}(\eta,\xi,h)
\;=\; \sum_{i=1}^{r}\;
m^{(i)}_{\mathrm{out}}(\D,h)\,\Bigl[\;\underbrace{\Phi_i(\eta)}_{\text{pointwise in }x}\;\cdot\;
\bigl(\,m^{(i)}_{\xi}(\D,h)\,\xi\,\bigr)\Bigr],
\]
implemented as:
\begin{enumerate}[leftmargin=1.6em, itemsep=1pt]
\item $\Phi(\eta) = \texttt{Dense}(256 \to r)\,E(\eta)$ --- a \textbf{zero-initialized}
      linear projection of the shared trunk features to one spatial gain field
      per branch (this is what makes the whole block inert at init);
\item $\xi$-branching in $k$-space: $\widehat{\xi}_i = m^{(i)}_{\xi}(k,h)\,\hat\xi(k)$,
      one filtered copy of the \emph{physical} $\xi$ per branch;
\item a \emph{pointwise} product $\Phi_i(\eta)(x)\cdot \xi_i(x)$ --- deliberately
      diagonal over branches (no outer product between $\eta$-channels and
      $\xi$-branches), so cost and memory are linear in $r$;
\item per-branch output multipliers $m^{(i)}_{\mathrm{out}}$, then a sum over branches.
\end{enumerate}
The full model is the baseline plus the block sum, rescaled to normalized units:
\[
\boxed{\;
\Gop_\theta(\eta;h)\,\xi
\;=\; \Gop_0(h)\,\xi
\;+\; \sum_{b=1}^{8} \mathrm{Block}_b\bigl(E(\eta),\ \xi,\ h\bigr)\;}
\]
Linearity in $\xi$ (P1) holds \emph{by construction}: $\xi$ enters every branch
exactly once, through a linear filter, a multiplication by an $\eta$-dependent
field, and another linear filter. For fixed $(\eta,h)$ the model \emph{is} a
linear operator in $\xi$ --- a sum of $8 \times 256 = 2048$
multiplier--gain--multiplier sandwiches. Each block can represent
$\Gop_1$-shaped terms exactly when $\Phi$ recovers an affine function of $\eta$,
and higher-order $\Gop_j$ effectively, since the trunk makes $\Phi_i(\eta)$ an
arbitrary (learned, pointwise) nonlinear functional of the seven $\eta$-channels.

\subsection{Expressivity notes and deliberate trade-offs}
\label{sec:expressivity}

\begin{itemize}[leftmargin=1.6em, itemsep=2pt]
\item \textbf{Zero-phase multipliers.} Real $m(k,h)$ on rfft modes are
      even, zero-phase convolution kernels. Odd differential action on $\xi$ is
      therefore not exactly representable: of the analytic
      $\Gop_1 = -\Gop_0\,\eta\,\Gop_0 - \partial_x\,\eta\,\partial_x$, the first
      term and the $\eta\,\partial_x^2$ part of the second are exactly reachable
      ($\Gop_0$ and $k^2$ are real multipliers, with $\eta,\partial_x^2\eta$
      available to $\Phi$), while the $\eta_x\,\partial_x$ part can only be
      approximated through combinations (the $\eta$-side does have odd operators
      via $\partial_x\eta$ and $\Hil\eta$). This is a deliberate economy ---
      complex multipliers double the parameter count of the dominant tensors ---
      and at the achieved single-step error ($\sim 10^{-3}$ relative $H^1$) it
      has not been the binding constraint.
\item \textbf{No enforced self-adjointness.} The true $\Gop$ is self-adjoint and
      positive semi-definite; the baseline run learns $m_\xi \ne m_{\mathrm{out}}$
      independently, so symmetry is only approximate. A
      \texttt{tie\_xi\_out\_mult} flag exists that sets
      $m_\xi = m_{\mathrm{out}}$ per branch, making each branch manifestly
      self-adjoint (see \S\ref{sec:parked}).
\item \textbf{Exact flat-surface limit only at init.} With biases present,
      $\Phi(0) \ne 0$ after training, so $\Gop_\theta(0;h) \ne \Gop_0(h)$
      exactly (it is close in practice). A \texttt{phi\_bias\_free} flag removes
      all trunk/projection biases so the residual vanishes identically at
      $\eta=0$ for \emph{all} weights (\S\ref{sec:parked}).
\item \textbf{Global support in $x$.} Although $\Phi$ is pointwise, the
      multipliers are dense in $k$, so the learned operator has full nonlocal
      reach --- consistent with the true DNO, which is nonlocal.
\end{itemize}

\section{Hyperparameters and parameter budget}

\begin{table}[h]
\centering
\small
\begin{tabular}{llll}
\toprule
field & value & meaning \\
\midrule
\texttt{width}      & 512 & trunk feature width is \texttt{width}/2 = 256 \\
\texttt{n\_blocks}  & 8   & Craig--Sulem blocks summed \\
\texttt{latent}     & 256 & branches $r$ per block \\
\texttt{cs\_n\_polys} & 3 & $\eta, \eta^2, \eta^3$ features \\
\texttt{cs\_use\_*\_deriv / hilbert} & all true & $\partial_x, \partial_x^2, |\D|^{1/2}, \Hil$ features \\
\texttt{cs\_mult\_hidden} & 128 & multiplier MLP hidden width \\
\texttt{h\_clip\_max} & 5.0 & depth saturation \\
\texttt{modes} & 64 & \emph{unused} (CLI compatibility with FNO/SpectralDNO) \\
\bottomrule
\end{tabular}
\caption{Baseline architecture hyperparameters
(\texttt{cs\_dno\_w512b8\_l256\_v3\_20260605\_032608}).}
\end{table}

\begin{table}[h]
\centering
\small
\begin{tabular}{lrr}
\toprule
component & params & share \\
\midrule
$\eta$ trunk (\texttt{eta\_feat\_proj} $7{\to}256$ + \texttt{eta\_feat\_mix} $256{\to}256$) & 67{,}840 & 6.0\% \\
$\Phi$ projections ($8 \times$ \texttt{Dense} $256{\to}256$, zero-init) & 526{,}336 & 46.5\% \\
multiplier MLPs ($8$ blocks $\times$ $\{m_\xi, m_{\mathrm{out}}\}$ $\times$ ($4{\to}128{\to}256$)) & 538{,}624 & 47.6\% \\
\midrule
total & 1{,}132{,}800 & \\
\bottomrule
\end{tabular}
\caption{Exact parameter budget (matches \texttt{param\_count} in the run
config). For comparison, the FNO baseline \texttt{fno\_w128b6\_v3\_hclip5} has
13.1M parameters.}
\end{table}

The model is small because it learns only scalar gain fields and 1-D multiplier
curves on top of hard-coded structure --- there are no $512\times512\times64$
complex spectral weight tensors anywhere.

\section{Training protocol}

\textbf{Data.} \texttt{combined\_dataset\_v3.npz}: 5{,}993{,}989 snapshots
$(\eta, \xi, \Gop\xi, h)$ at $n_x = 1024$, harvested from f64 solver rollouts
across ten families --- Tanaka solitons g0/g1 (2.0M), Benjamin--Feir g0/g1/modal
(2.03M), Stokes deep/finite (1.0M), linear shallow sweeps (0.5M), random seas
deep/finite (0.46M). Labels are the order-6/pad-8 operator \emph{lowpassed to
the $k\le128$ subspace} (the pipeline's invariant subspace). Split 80/10/10 by a
seeded global permutation: 4{,}795{,}193 train / 599{,}398 val. The v4 dataset in
the current retrain appends a shallow--steep wavetrain family ($+480$k, source
11) and changes nothing else; see \texttt{v4\_dataset\_report.pdf}.

\textbf{Loss.} Relative Sobolev $H^1$ in the spectral domain, on normalized
targets:
\[
\mathcal{L} \;=\; \operatorname{mean}_{\text{batch}}
\frac{\bigl\|\,w \odot (\widehat{\mathrm{pred}} - \widehat{\mathrm{tgt}})\,\bigr\|_2}
     {\bigl\|\,w \odot \widehat{\mathrm{tgt}}\,\bigr\|_2},
\qquad w(k) = \sqrt{1 + k^2},
\]
i.e.\ an $L^2$-plus-gradient norm that up-weights exactly the high-$k$ content a
rollout differentiates; relative (per-sample) so soliton rows and near-flat rows
contribute comparably.

\textbf{Optimization.} AdamW, base lr $2\times10^{-4}$ with single cosine decay
over the full run, weight decay $10^{-4}$, global batch 256, 40 epochs, f32
training throughout. The baseline trained on one GPU; the trainer shards the
global batch across devices via \texttt{shard\_map} with per-step gradient
all-reduce (the v4 retrain uses both GPUs this way --- optimization-identical to
single-device). Best-val checkpoint is kept: \textbf{val $= 9.28\times10^{-4}$}
relative $H^1$, vs $3.69\times10^{-3}$ for the 13.1M-parameter FNO baseline on
the same data and loss (80 epochs) --- $4\times$ lower error with $11.6\times$
fewer parameters at half the epochs.

\section{Deployment in rollouts}

\texttt{model\_rollout.py} reconstructs the module from \texttt{config.json}
and substitutes it for the series evaluation inside the same time integrator
used to generate truth, with the same $0.25$ dealias filter applied to the RHS
and post-step states --- the surrogate rollout lives in the same invariant
subspace it was trained on. The network always evaluates in f32; with
\texttt{--f64\_harness} (the honest configuration --- see
\texttt{v4\_dataset\_report.pdf}) all surrounding integrator state and spectral
algebra run in f64. Depth enters the surrogate as $\log h$ exactly as in
training.

\section{Empirical status (baseline, f64 harness, truth-gated)}

\begin{table}[h]
\centering
\small
\begin{tabular}{lcc}
\toprule
regime & divergence & median rel-$L^2(\eta)$ at $t{=}200$ \\
\midrule
bf\_g0      & 0\%    & 0.0044 \\
bf\_g1      & 0\%    & 0.0067 \\
bf\_modal   & 6.7\%  & 0.0055 \\
tanaka\_g0  & 6.2\%  & 0.0420 \\
tanaka\_g1  & 18.8\% & 0.0443 \\
linear (shallow-steep sweep) & 27\% & 0.084 \\
\bottomrule
\end{tabular}
\caption{Long-rollout accuracy of the v3 baseline. The single systematic
weakness is shallow water ($h \lesssim 0.1$) at moderate steepness
($a/h \gtrsim 0.15$) on non-soliton states --- a training-coverage gap, not an
architecture failure (single-step error on exact solitons is
$\sim 2\times10^{-4}$ even where rollouts die). The v4 retrain targets exactly
this with the shallow--steep family.}
\end{table}

\section{Parked architectural variants}
\label{sec:parked}

Three flags exist in the module, are plumbed through training and rollout, and
are \textbf{off} in the baseline; the program that motivated them was closed
when the f32 FFT noise floor was identified as the real obstacle:

\begin{itemize}[leftmargin=1.6em, itemsep=2pt]
\item \texttt{use\_g1\_baseline} --- add the closed-form
      $\Gop_1(\eta;h)\xi = -\Gop_0(\eta\,\Gop_0\xi) - \partial_x(\eta\,\partial_x\xi)$
      to the analytic baseline so learned blocks model only $\Gop_{2+}$.
      Requires \texttt{g1\_k\_cut} (default 128): the doubly $k$-amplified f32
      FFT chain produces rounding noise above $k\approx128$ that exceeds the
      genuine $\Gop_1$ signal there, which the $H^1$ loss then up-weights into
      an unlearnable $\sim3\times10^{-3}$ floor --- the masked variant zeroes
      the term's output modes at $k \ge$ cut.
\item \texttt{tie\_xi\_out\_mult} --- share one multiplier as both $m_\xi$ and
      $m_{\mathrm{out}}$ per branch, making every branch (hence every block)
      self-adjoint in $\xi$, matching the true operator's symmetry.
\item \texttt{phi\_bias\_free} --- remove all biases from the $\eta$ trunk and
      $\Phi$ projections so every learned term vanishes identically at
      $\eta = 0$: the flat-surface limit $\Gop_\theta(0;h) = \Gop_0(h)$ becomes
      exact for all trained weights, not just at init.
\end{itemize}

\end{document}
